% !TEX root = main_memo.tex

\chapter{Finite generation at double successors}\label{sectionDoubleSucc}

In view of \cref{FGconsequences}, it remains to show that $\FG(<\alpha+2)$ implies $\FG(\alpha+2)$ in order to complete the proof by induction of \cref{PreciseStructureThm}.
To do so, writing $\alpha=\lambda+n$ with $\lambda$ limit and $n\in\N$, we can can rely on a decomposition in centered functions, which are either in $\sC_{<\lambda}$ or in $\centered{\alpha+2}$, up to equivalence.
The situation is more complex than for successor of limits, however we can actually assume that they are all in $\centered{\alpha+2}$ thanks to \cref{ConsequencesGeneralStructureThm}~\cref{ConsequencesGeneralStructureThm2} (see \cref{lem:gobblingLessThanLambda}).
For this general case, we show that partitions in centered functions can be made particularly well-behaved. We then analyze two gradually more complex situations before we define solvable functions. We finally show that solvable functions are equivalent to finite gluings of our generators and how to reduce the general case to that of solvable functions.


Note that this strategy departs from the one we used at successors of limit (\cref{FGatsuccessoroflimit}) which relied on two structural properties of continuous reducibility specific to these levels. %at successor of limit levels.
First, by \cref{cor:CenteredSucessor} there are up to equivalence only two centered functions: the minimum and the maximum among simple functions, which makes them comparable.
Second, the only generator at this level that requires the wedge operation,
namely $\fctwedge(\Maximalfct{\lambda}\mid\Minimalfct{\lambda+1})$, is actually of infinite $\CB$-degree.

 

\section{Fine partitions in centered functions}\label{DoubleSucc_domain}

We introduce some terminology and notations for partitions in centered functions.

Given topological spaces $A$ and $B$, a \emph{$c$-partition}\index[IdxNotion]{partition!$c$-partition} of a function $f:A\to B$ in $\sC$ is a clopen partition $\partit\subseteq \mathbf{\Delta}^0_1(A)$ of $A$
such that for all $P\in \partit$ the restriction $f\restr{P}$ is centered with \cocenter{} denoted by $y_P$.
\index[IdxSymb]{P@{$\partit$}}\index[IdxSymb]{y@{$y_P$}}

Let $\partit$ be a $c$-partition for $f$. We let $Y_\partit=\set{y_P}[ P\in \partit]\subseteq B$ be the (countable) set of all \cocenter{}s of the functions $f\restr{P}$ for $P\in \partit$.
For a centered function $g$ and $y\in Y_\partit$, let $\partit_{g,y}=\set{P\in \partit}[f\restr{P}\equiv g \text{ and } y_P=y]$ and $f_{(g,y)}=\bigsqcup_{P\in \partit_{g,y}}f\restr{P}$.
\index[IdxSymb]{Y@{$Y_\partit$}}\index[IdxSymb]{P@{$\partit_{g,y}$}}\index[IdxSymb]{f@{$f_{(g,y)}$}}

\subsection{Dissolving lumps}

Let $f\in \sC_{\alpha}$ for $\alpha=\lambda+n$ with $\lambda$ limit and $n\in \N$. We define $\omegaregular{\alpha}=\set{\Maximalfct{\lambda}}\cup \set{\omega h }[h\in\centered{\alpha}]$. 
\index[IdxSymb]{W@{$\omegaregular{\alpha}$}}
We say that $f$ is \emph{$\omegaregular{}$-regular}\index[IdxNotion]{function!$\omegaregular{}$-regular at a point}\index[IdxNotion]{w@{$\omegaregular{}$-regular function}} at $y\in B$ if for all $h\in \omegaregular{\alpha}$ the set $\set{j\in \N}[ h\leq \ray{f}{y,j}]$ is either empty or infinite. 

\index[IdxNotion]{lump}
Given a $c$-partition $\partit$ of $f$, a \emph{$\partit$-lump} is a pair $(g,y)$ with $y\in Y_\partit$ and $g$ centered such that $f_{(g,y)}$ is not $\omegaregular{}$-regular at $y$.
The \emph{rank} of the $\partit$-lump is the rank $\CB(g)$ of $g$, which is equal to $\CB(f_{(g,y)})$.

\medskip

\begin{remark}
A centered function\index[IdxNotion]{function!centered} $f$ is $\omegaregular{}$-regular at its \cocenter{} $y$, because for all $j,m$ we have $\ray{f}{y,j}\leq \gl_{m=i}^M\ray{f}{y,i}$ for some $M>m$ by \cref{Rigidityofthecocenter}. Hence if $\omega h\leq \ray{f}{y,j}$ for some $h\in \centered{\alpha}$, then $\omega h\leq \ray{f}{y,m}$ for arbitrarily large $m$ by \cref{Intertwinereductionsforomegacentered}. {Similarly for $\Maximalfct{\lambda}=\gl_n \Maximalfct{\beta_n}$ with $(\beta_n)_n$ cofinal in $\lambda$. } Therefore $f_{(g,y)}$ is not centered when $(g,y)$ is a lump.
In particular, it follows that if $(g,y)$ is a lump, then $\partit_{g,y}$ is infinite. {Because if $\partit_{g,y}$ is finite, then $f_{(g,y)}\equiv g$ is centered, hence $\omegaregular{}$-regular. To see this assume that $g\equiv \pgl_n g_n$ for a monotone sequence $(g_n)_n$ by \cref{CenteredasPgluing}. For all $P\in \partit_{g,y}$, the sequence of rays of $f\restr{P}$ along $y$ is reducible by pieces to $(g_n)_n$ by \cref{Rigidityofthecocenter}. Now each ray of $f_{(g,y)}$ at $y$ is the finite disjoint union of corresponding rays of $f\restr{P}$, for $P\in \partit_{g,y}$. So the rays of $f_{(g,y)}$ at $y$ are reducible by pieces to $(g_n)_n$. Therefore $g\leq f_{(g,y)}\leq \pgl_n \ray{(f_{(g,y)})}{y,n} \leq \pgl_n g_n$ by \cref{Pgluingasupperbound}.}
% \ynote{Proof that it is centered. Do we need it? Or show directly that it is $\omegacentered{}$-regular: To see this, note that the rays of $f_{(g,y)}$ are equal to $\bigsqcup_{P\in \partit_{g,y}}\ray{f\restr{P}}{y, n}$, $n\in \N$. Hence if $\partit_{g,y}$ is finite and $\omega h\leq \ray{f_{(g,y)}}{y,n}$ for some $h\in \centered{\alpha}$, then $\omega h\leq \ray{f\restr{P}}{y,n}$ for some $P\in \partit_{g,y}$ by \cref{Intertwinereductionsforomegacentered}.}
\end{remark}

Recall that a partition $\partit'$ is \emph{finer}\index[IdxNotion]{partition!finer} than a partition $\partit$ if every $P'\in\partit'$ is included in some $P\in\partit$.
A $c$-partition presenting a lump can be refined to \emph{dissolve} that lump.

\begin{lemma}\label{lemma:RefiningBy1}
Let $\alpha<\omega_1$ and assume $\FG(<\alpha)$. Let $f\in \sC_\alpha$ and $\partit$ a $c$-partition for $f$.

If $(g,y)$ is a $\partit$-lump of rank $\beta\leq\alpha$, then there exists a finer $c$-partition $\partit'$ such that $(g,y)$ is not a $\partit'$-lump,
 $\partit\setminus\partit_{(g,y)}\subseteq\partit'$, and every $\partit'$-lump is either a $\partit$-lump or has rank $<\beta$.
\end{lemma}
\begin{proof}
Suppose that $(g,y)$ is a $\partit$-lump of rank $\beta$, namely that $h=f_{(g,y)}$ is not $\omegaregular{}$-regular.
Note that $\CB(h)=\CB(g)=\beta$ by \cref{CBrankofclopenunion}.
There exists $w\in \omegaregular{\beta}$ such that $J_w=\{j\in\N\mid w\leq \ray{h}{j}\}$ is finite but non-empty.
Since $\omegaregular{\beta}$ is finite, there exists $J\in \N$ such that for all $w\in \omegaregular{\beta}$ either $J_w\subseteq J$ or $J_w$ is infinite. Let $U=N_{y\restr{J+1}}$ and note that $h\corestr{U}$ is $\omegaregular{}$-regular.

Now for all $P\in \partit_{(g,y)}$, let $P'=P\cap f^{-1}(U)$ and $A_P=P\setminus P'$. Note that while $f\restr{P'}\equiv g$ by centeredness and we have $\CB(f\restr{A_P})<\CB(g)=\beta\leq\alpha$ by \cref{CenteredasPgluing}.
So by $\FG(<\alpha)$ we can use \cref{FGconsequences} to get $\partit_P$ a $c$-partition of $f\restr{A_P}$ with functions of $\CB$-rank $<\beta$.
Our new $c$-partition for $f$ is given by 
\[
\partit'=(\partit\setminus \partit_{g,y})\cup \{P' \mid P\in \partit_{g,y}\} \cup \bigcup \{\partit_P \mid P\in \partit_{g,y}\}\,.
\]
Note that $\partit'_{g,y}=\{P'\mid P\in \partit_{g,y}\}$ and $\bigcup \partit'_{g,y}=f^{-1}(U)\cap \bigcup \partit_{g,y}$, so $(g,y)$ is not a $\partit'$-lump. %by construction.
Moreover for every $Q\in \partit'\setminus \partit$ if $Q\notin \partit'_{g,y}$ then $\CB(f\restr{Q})<\CB(g)=\beta$ holds, which proves the statement.
\end{proof}


\subsection{Gobbling up small functions}

Writing $\alpha=\lambda+n$ with $\lambda$ limit and $n$ finite, note that a $c$-partition of $f\in\sC_\alpha$ may contain a clopen $P$ such that $\CB(f\restr{P})<\lambda$.
When $\alpha$ is a double successor, that is when $n>1$, the parts of $\CB$-rank bigger than $\lambda+1$ can gobble up those of rank $<\lambda$.

\begin{lemma}\label{lem:gobblingLessThanLambda}
Let $\lambda<\omega_1$ be limit and $f\in\sC$.
Assume that $f=f_0\sqcup f_1$ with $f_0$ centered, $\pgl \Maximalfct{\lambda}\leq f_0$ and $f_1\leq \Maximalfct{\lambda}$. Then $f$ is centered and $f\equiv f_0$.
\end{lemma}
\begin{proof}
Let $x$ be a center for $f_0$ and let $U\ni x$ be clopen in $\dom (f_0)$. Setting $f_U=f\restr{U}$, we have $f_{U}\leq f_0$ and we show that $f\leq f_{U}$.

Since $x$ is a center for $f_0$, we have $f_0\leq f_{U}$ and in particular $\pgl \Maximalfct{\lambda}\leq f_{U}$ as witnessed by some reduction $(\sigma,\tau)$.
Let $\tilde{y}=f\sigma(\iw{0})$ and $y=f(x)$. If $\tilde{y}\neq y$, choose a clopen neighborhood $V\subseteq \im (f)$ of $\tilde{y}$ not containing $y$.
Otherwise, set $V=\tau^{-1}(N_{(1)})$, notice that $V$ is a clopen set with $y\notin V$, and such that $ \Maximalfct{\lambda}\leq f_{U}\corestr{V}$ as witnessed by restricting $(\sigma,\tau)$.
Now let $W\ni y$ be clopen and disjoint from $V$. By centeredness of $f_0$, we have $f_0\leq f_{U}\corestr{W}$. Moreover $f_1\leq\Maximalfct{\lambda}\leq f_{U}\corestr{V}$. Hence by \cref{Gluingasupperbound,Gluingaslowerbound}, we get
\[
f\leq f_0 \glbin f_1 \leq  f_{U}\corestr{W} \glbin  f_{U}\corestr{V} \leq f_{U}.\qedhere
\]
\end{proof}

\subsection{Fine $c$-partitions.}\label{DefFine}
We now define the partitions with all the required properties. 
\begin{definition}
Let $f\in \sC_{\lambda+n+1}$ for $\lambda$ limit and $n\in \N$. We say that a $c$-partition $\partit$ of $f$ is \emph{fine} if there are no $\partit$-lumps and $\CB(f\restr{P})>\lambda$ for all $P\in \partit$.\index[IdxNotion]{fine $c$-partition}\index[IdxNotion]{partition!fine $c$-partition}

\end{definition}

\medskip

We now show that $\FG(<\alpha)$ yields the existence of fine $c$-partitions when $\alpha$ is a double successor.
To do so, we take any $c$-partition and start by dissolving all lumps of large rank before gobbling up any small pieces that it may contain.

\index[IdxNotion]{finitely generated set of functions}
\begin{proposition}\label{ExistenceFinePartitions}
Let $\alpha=\lambda+n+2$ with $\lambda<\omega_1$ limit, $n\in\N$, and assume $\FG(<\alpha)$.
Then every function in $\sC_\alpha$ admits a fine $c$-partition.
\end{proposition}
\begin{proof}
Take $f\in \sC_\alpha$, and start by defining a sequence $(\partit_i)_{i\leq n+2}$ of $c$-partitions of $f$ such that all $\partit_i$-lumps have rank $\leq\lambda+n+2-i$.
By $\FG(<\alpha)$, $f$ admits a  $c$-partition $\partit_0$ by \cref{FGconsequences}.

Suppose by induction that $\partit_i$ is defined for some $i\leq n+2$ and fix an enumeration $(g_j,y_j)_{j\in J}$ for $J$ an initial segment of $\N$ of all $\partit_i$-lumps of rank $\lambda+n+2-i$.
We define by induction a sequence $(\partit'_j)_{j\in J}$ of $c$-partitions, starting with $\partit'_{-1}=\partit_i$. For all $j\in J$, let $\partit'_{j+1}$ be a $c$-partition refining $\partit'_{j}$ and dissolving the $\partit_i$-lump $(g_j,y_j)$ by an application of \cref{lemma:RefiningBy1}.
The limit inferior $\partit_{i+1}=\bigcup_{k\in J} \bigcap_{k\leq j \in J} \partit'_j$ is the desired $c$-partition.
Note that all $\partit_{n+2}$-lumps have rank $<\lambda$. Moreover since $\CB(f)\geq \lambda+2$ we have $\pgl \Maximalfct{\lambda}\leq f$ by \cref{ConsequencesGeneralStructureThm}~\cref{ConsequencesGeneralStructureThm2}. Therefore $\pgl \Maximalfct{\lambda}\leq f\restr{P}$ for some $P\in \partit_{n+2}$ by \cref{Centerinvariance}~\ref{Centerinvariance3}.
So letting $\tilde{\partit}=\{Q\in \partit_{n+2} \mid f\restr{Q}\leq \Maximalfct{\lambda}\}$ and $R=\bigcup \tilde{\partit}$,
then $\CB(f\restr{P\cup R})\geq \CB(\pgl \Maximalfct{\lambda})>\lambda$ and by \cref{lem:gobblingLessThanLambda}
$f\restr{P\sqcup R}$ is centered and $f\restr{P\sqcup R}\equiv f\restr{P}$ .
Therefore the $c$-partition $\partit$ obtained from $\partit_{n+2}$ by replacing $P$ and elements of $\tilde{\partit}$ by $P\cup R$ is a fine $c$-partition.
\end{proof}



\section{Pseudo-centered functions}\label{subsec:pseudocentered}

Assuming $\FG(<\alpha+2)$ any function in $\sC_{\alpha+2}$ admits a fine $c$-partition $\partit$ by \cref{ExistenceFinePartitions}.
Hence every $P\in \partit$ is associated with both a centered function in $\centered{\alpha+2}$ and a point $y_P$ in the codomain of $f$.
{The difficulty in proving that $f$ is equivalent to a finite gluing of generators in $\generator{\alpha+2}$
resides in the potentially complex combinatorial and topological structure of the combination of the mappings $\partit\to \centered{\alpha+2}$ and $\partit\to \im f$. }
Before tackling the general case, we deal with ubiquitous situations that are easier to handle.
The first one consists of a disjoint union of the same centered function with only one \cocenter{},
in other words (using the notations of \cref{DoubleSucc_domain}) a function $f=f_{(g,y)}$ for some centered $g$ and $y\in\im(f)$.

\index[IdxNotion]{pseudo-centered function}
\begin{definition}
We say that $f\in\sC$ together with a fine $c$-partition is \emph{pseudo-centered} at $y$ if $Y_\partit$ is a singleton $\set{y}$ and for all $P,P'\in \partit$ we have $f\restr{P}\equiv f\restr{P'}$.
\end{definition}

The following result generalizes \cref{simplefunctionslambda+1damuddafuckaz} to functions whose $\CB$-rank is a double successor.
The strategy of proof is an elaboration on the successor of limit case.

\index[IdxNotion]{Precise Structure!pseudo-centered functions}
\index[IdxNotion]{Vertical Theorem}
\index[IdxNotion]{generators}
\index[IdxNotion]{finitely generated set of functions}
\begin{theorem}[Vertical theorem]\label{VerticalTheorem}
Let $\alpha<\omega_1$ and assume $\FG(\leq\alpha+1)$.
Let $f:A\to B$ in $\sC_{\alpha+2}$ be pseudo-centered at $y$.
There exist $g\in\centered{\alpha+2}$ and $H\subseteq \omegaregular{\alpha+1}$ such that
for all clopen neighborhood $U$ of $y$ there is a clopen set $W\subseteq U$ and a clopen partition $A=A^0\sqcup A^1$ such that 
\begin{enumerate}
\item $y\notin W$ and $f\restr{A^0}\leq \gl H \leq f\corestr{W}$,
\item for all clopen set $V\ni y$, $f\restr{A^1}\leq g \leq f\corestr{V}$ (in fact  $g \leq f\restr{A^1}\corestr{V}$),
\item $\gl H\leq g$, so in particular $f\leq g\gl g$.
\end{enumerate}
\end{theorem}
\begin{proof}
Fix a fine $c$-partition $\partit$ such that $Y_\partit=\set{y}$ and some centered $\hat{f}$ with $f\restr{P}\equiv \hat{f}$ for all $P\in \partit$.
In this proof, all rays are taken with respect to $y$, so that $\ray{f}{j}$ denotes the rays of $f$ at $y$. Note that $\CB(\ray{f}{j})\leq\alpha+1$.
Since $\alpha+2=\CB(f)=\CB(\hat{f})\neq\lambda+1$ for any limit $\lambda\leq\alpha$, we have $\hat{f}>\Minimalfct{\lambda+1}$,
so $\hat{f}\equiv \pgl G$ for some non-empty set $G$ of $\centered{\alpha+1}$ by \cref{FGconsequences} and so $\pgl G\in\centered{\alpha+2}$. 
Note that for all clopen set $V\ni y$, $\pgl G \leq f\corestr{V}$, because for any $P\in \partit$ we have $\pgl G\leq f\restr{P} \leq f\restr{P}\corestr{V}$.

If for all $j\in\N$ we have $\ray{f}{j}\leq\FinGl{G}$, then $f\leq\pgl_j\ray{f}{j}\leq \pgl G\equiv g$ by \cref{Pgluingasupperbound}.
Observe that this happens automatically if $\partit$ is finite.
In this case, we can set $w=\emptyset$, $W=\emptyset$, $A^0=\emptyset$ and $A^1=A$. 

\smallskip

Otherwise, $\partit$ must be infinite and we consider the set 
\[
H=\set{h\in \omegaregular{\alpha+1}}[h\nleq\FinGl{G} \text{ and } \exists j\in\N \; h \leq \ray{f}{j}]
\]
%$H$ of functions $\tilde{w}$ in \ynote*{maybe take $H\subseteq \centered{\alpha+1}$ and define $w=\omega H$, rework with the final definition of $\omegacentered{\alpha+1}$, here any way $\Maximalfct{\lambda}\leq G$.}{$\omegacentered{\alpha+1}$ such that $\tilde{w}\not\leq\FinGl{G}$ and there exists $j\in\N$ with $\tilde{w}\leq \ray{f}{j}$.}
%We let $w=\gl H$.
We let $w=\gl H$ and we claim that for all clopen $U\ni y$ we have $w\leq f\corestr{W}$ for some clopen $W$ with $y\notin W\subseteq U$.

To see this, let $h\in H$. Since $(f,y)$ is not a $\partit$-lump, we must have $h\leq\ray{f}{j}$ for infinitely many $j\in\N$.
Suppose that $U\ni y$ is clopen, and therefore contains all but finitely many rays $\ray{B}{j}$.
As $H$ is a subset of the finite set $\omegaregular{\alpha+1}=\set{\Maximalfct{\lambda}}\cup\omega \set{\centered{\alpha+1}}$ by \cref{BasicsOnGenerators}, there is a finite set $J\subseteq\N$ such that for all $h\in H$ there exists $j\in J$ with
$h\leq \ray{f}{j}=f\corestr{\ray{B}{j}}$ and $\bigcup_{j\in J}\ray{B}{j}\subseteq U$.
Setting $W=\bigcup_{j\in J}\ray{B}{j}$ and $w\leq f\corestr{W}$ by \cref{Intertwinereductionsforomegacentered}.%\ynote*{check and adapt}{it follows that $w=\gl H\leq f\corestr{W}$ by \cref{BasicsOnGenerators} again}.

\smallskip

Next we show that $w\leq \pgl G$.  Note that $\ray{f}{j}=\bigsqcup_{P\in \partit} \ray{f\restr{P}}{j}$ and $\ray{f\restr{P}}{j} \leq \FinGl G$ by \cref{Rigidityofthecocenter} and so $\ray{f}{j}\leq \omega G$ for all $j$.
Hence by definition of $H$, {$h\leq \omega G$ for $h\in H$ and thus $w\leq \omega G$}. Since $\omega G\leq \pgl G$ by \cref{GluinglowerthanPgluing}, it follows that  $w\leq \pgl G$.

It remains to find a clopen partition $A=A^0\sqcup A^1$ such that $f\restr{A^0}\leq w$ and $f\restr{A^1}\leq g$.
Fix an enumeration $(A_i)_{i\in\N}$ of $\partit$, let $f_i=f\restr{A_i}$, and for all $j\in\N$ set $f^{[j]}=\bigsqcup_{i>j} \ray{f_i}{j}$. We prove the following:

\begin{figure}
%\includegraphics[width=0.3\linewidth]{verticalthm_diagonalsplit.png}
\centering
\index[IdxFig]{Proof of the Vertical Theorem}
\input{verticalthm_diagonalsplit_small.tex}
\caption{Proof of the Vertical Theorem.}
\end{figure}
\begin{claim}
For all $j\in \N$, there exists a clopen partition $(A^0_j, A^1_j)$ of $\dom(f^{[j]})$
such that:  $f^{[j]}\restr{A^0_j}\leq w$ and $f^{[j]}\restr{A^1_j}\leq \FinGl{G}$.
\end{claim}
\begin{subproof}
Let $j\in \N$, we show that $f^{[j]}\leq w \glbin h$ for some $h\in\FinGl{G}$.
By \cref{Gluingasupperbound}, this grants us with the desired clopen partition $(A^0_j, A^1_j)$ of $\dom(f^{[j]})$. 

Since $\CB(f^{[j]})< \alpha+2$, either $\CB(f^{[j]})< \lambda$ or by $\FG(\leq\alpha+1)$ it is equivalent to a finite gluing of functions in $\generator{\alpha+1}$.
In the former case $f^{[j]}\leq G$ by \cref{JSLgeneralstructure} and we are done.
In the latter case, we can write $f^{[j]}=\gl_{i\leq k}g_i$ for some $k\in\N$ and $g_i\in\generator{\alpha+1}$.
So in order to prove the claim, as $\generator{\alpha+1}$ is finite, it suffices to show that for every generator $g\in \generator{\alpha+1}$ if $g\leq f^{[j]}$, then $g \leq w \glbin h$ for some $h\in\FinGl{G}$.
So let $g\in \generator{\alpha+1}$ with $g\leq f^{[j]}$, we distinguish three cases according to the definition of $ \generator{\alpha+1}$:
\begin{enumerate}
\item Suppose that $g\in \centered{\alpha+1}$, so $g$ is centered. Since $f^{[j]}=\bigsqcup_{i>j} \ray{f_i}{j}\leq \gl_{i>j}\ray{f_i}{j}$ by \cref{Gluingasupperbound}
and $\ray{f_i}{j}\leq \FinGl G$ for all $i,j\in\N$ by \cref{Rigidityofthecocenter}, we actually have $g\leq G$ by \cref{Centerinvariance}.

\item Suppose that $g\in\omegaregular{\alpha+1}$. By definition of $H$, we have either $g\leq\FinGl{G}$ or, since $g\leq f^{[j]}\leq \ray{f}{j}$, $g\in H$ and so $g\leq w$.

\item Finally\footnote{Here we see the specificity of double successors!},
it can be that $g$ is a wedge generator, so use \cref{BasicsOnGenerators} to get $F_0,F_1\subseteq \centered{\alpha+1}$ such that $F_0,\omega F_1\leq g$ and $g\leq(\gl F_0)\glbin \omega F_1$

Since $F_0\leq f^{[j]}$, we have $F_0\leq G$ by the first case.
Since $\omega F_1\leq f^{[j]}$, either $\omega F_1 \leq\FinGl{G}$ or $\omega\set{F_1}\subseteq H$ by definition of $H$, and so $\omega F_1\leq \gl H= w$.
This shows that for some $h\in \FinGl{G}$ we have
\[
g\leq\gl \omega F_1 \glbin F_0  \leq w \glbin h, 
\]
as desired.\qedhere
\end{enumerate}
\end{subproof}


We let $A^0=\bigcup_{j\in\N} A^0_j$ and $A^1=A\setminus A^0$;
$A^0$ is open as union of the sets $A^0_j$ which are open, we prove $A^1$ is also open. 
Whenever $j\geq i$, we have $A_i\cap \dom(f^{[j]})=\emptyset$ and so for all $i\in \N$ the set $A_i\cap A^0=A_i \cap \bigcup_{j<i}A^0_j$ is clopen, hence so is $A_i\setminus A^0$.
Therefore $A^1=A\setminus A^0=\bigcup_i A_i\setminus A^0$ is open.

On the one hand, 
\[
f\restr{A^0}\leq \gl_{j} f\restr{A^0_j} \leq \omega w \equiv w.
\] 

On the other hand, for all $j\in\N$ we have by \cref{Gluingasupperbound}
\[
\ray{f\restr{A^1}}{j}=f\restr{A^1_j} \sqcup  \bigsqcup_{i\leq j}\ray{f_i}{j}
\leq f\restr{A^1_j}\glbin \gl_{i\leq j} \ray{f_i}{j} \leq \FinGl G,
\]
where the last reduction follows from the claim and the fact that $\ray{f_i}{j}\leq \FinGl G$ for all $i,j\in\N$. Therefore $(\ray{f\restr{A^1}}{j})_j$ is piecewise reducible to $\iw{\gl G}$, so by \cref{Pgluingasupperbound} $f\restr{A^1}\leq\pgl G$, which concludes the proof. 
\end{proof}

\section{Strongly solvable functions}
For a fine $c$-partition of a function $f$, we denote by $\partitn{y}=\set{P\in\partit}[y_P\neq y]$\index[IdxSymb]{P@{$\partitn{y}$}} the members of the partition whose \cocenter{}s are distinct from $y$.
Now that we dealt with the case of a unique \cocenter{} $y$, we next analyze the situation where
cocenters $\{y_P\mid P\in\partitn{y}\}$ converge to $y$ and the $\set{f\restr{P}}[P\in\partitn{y}]$ also enjoys a nice combinatorial property.
\begin{definition}
We say that the function $f:A\to B$ together with a fine $c$-partition $\partit$ is \emph{strongly solvable} at $y\in Y_\partit$ if for all clopen neighborhood $V$ of $y$\index[IdxNotion]{strongly solvable function}\index[IdxNotion]{function!strongly solvable}
\begin{enumerate}
\item the set $\{y_P\mid P\in\partit \text{ and } y_P\notin V\}$ is finite, and
\item for all $P\in\partitn{y}$ there exists $Q\in \partitn{y}$ with $y_{Q}\in V$ and $f\restr{P}\leq f\restr{Q}$.
\end{enumerate}
 \end{definition}

%Note that a centered function in $\sC$ with the trivial partition is strongly solvable at its \cocenter{}. 

\index[IdxNotion]{Precise Structure!strongly solvable functions}
\index[IdxNotion]{Diagonal Theorem}
\index[IdxNotion]{finitely generated set of functions}
\index[IdxNotion]{generators}
\begin{theorem}[Diagonal Theorem]\label{DiagonalTheorem}
Assume $\FG(\leq\alpha+1)$ for $\alpha<\omega_1$.
Let $f:A \to B$ in $\sC_{\alpha+2}$ be strongly solvable {at $y$}.
Then there exists $g\in \FinGl{\generator{\alpha+2}}$ such that $f\leq g\leq f\corestr{U}$ {for all clopen set $U\ni y$}. 
\end{theorem}

\begin{proof}
Let $\partit$ be a fine $c$-partition for $f$ and $\bar{y}\in Y_\partit$ witness that $f$ is strongly solvable {at $\bar{y}$}.
Let $\partit_M=\partit\setminus \partitn{\bar{y}}$ and note that if $\partitn{\bar{y}}$ is non-empty, then $Y'=\{y_P\mid P\in \partit\}\setminus\{\bar{y\}}$ is an infinite discrete subset of $B$ by strong solvability of $f$.
% in $U$.
We write $\alpha=\lambda+n$ with $\lambda$ limit and $n\in\N$ and distinguish two cases.
%There are two cases, we first treat the easier one and then the harder one.

\medskip

\textit{First case}: $\CB(f\restr{P})=\lambda+1$ for all $P\in \partit_M$. Since $\CB(f)=\alpha+2$, there exists $P\in\partit_D$ with $\CB(f\restr{P})=\alpha+2$ by \cref{CBrankofclopenunion}.
Choose a set of representatives $D\subseteq \centered{\alpha+2}$ for $\{f\restr{P}\mid P\in \partit_D\}$ by \cref{FGconsequences} and pick $h\in D$ with $\CB(h)=\alpha+2$.
We show that in that case we can set $g:=\omega D$.
By \cref{Maxfunctions,ConsequencesGeneralStructureThm} we have $f\restr{P}\leq\pgl\Maximalfct{\lambda}\leq h$ for all $P\in\partit_M$,
so using \cref{Gluingasupperbound,Gluingasupperbound_cor} we get:
\[
f \leq \gl_{P\in \partit} f\restr{P}\leq \gl_{P\in \partit_D} f\restr{P}\glbin\gl_{P\in \partit_M} f\restr{P}\leq \omega D\glbin \omega \pgl\Maximalfct{\lambda}\leq \omega D.
\]
Now fix any clopen set $U\ni \bar{y}$. Since $f$ is strongly solvable at $\bar{y}$, for all $g\in D$ the set $\set{y_P}[P\in\partit_D \text{ and }y_P\in U \text{ and }g\leq f\restr{P}]$ is infinite. So by \cref{Intertwinereductionsforomegacentered} we have $\omega D\leq f\corestr{U}$, which concludes the proof in this case.
%Note that $D$ is finite and for all $g\in D$ 
%$Y_D\cap V$ is infinite for all clopen neighbourhood $V\ni \bar{y}$, so by \cref{Intertwinereductionsforomegacentered} we have $\omega D\leq f\corestr{U}$.

\medskip

\textit{Second case}: $\CB(f\restr{P})> \lambda+1$ for some $P\in \partit_M$. 
Let $\partit'_M=\{P\in \partit_M\mid \CB(f\restr{P})> \lambda+1\}$ and choose a finite set of representatives
$M\subseteq \centered{\alpha+2}$ for $\{f\restr{P}\mid P\in \partit'_M\}$ by \cref{FGconsequences}.
For each $g\in M$, let $M_g\subseteq  \centered{\alpha+1}\cup \omega\{\alpha+1\}$ such that $g=\pgl M_g$ and let $A_g=\bigcup \{P\in \partit'_M\mid f\restr{P}\equiv g\}$. 

For all $g\in M$ the function $f\restr{A_g}$ is pseudo-centered, so by \cref{VerticalTheorem} there exists $H_g\subseteq \centered{\alpha+1}$,
a clopen $W_g\subseteq U$ and a clopen partition $A_g=A^0_g\sqcup A^1_g$ such that: first, $\bar{y}\notin W_g$ and $f\restr{A^0_g}\leq \omega H_g \leq f\restr{A^0_g}\corestr{W_g}$,
and second for all clopen $V\ni \bar{y}$ we have $f\restr{A^1_g}\leq g \leq f\restr{A^1_g}\corestr{V}$ (so in particular,  $f\restr{A^1_g}\equiv g$).


If $\partit_D$ is empty, set $D=\emptyset$. Otherwise, as in the previous case, let $D\subseteq \centered{\alpha+2}$ be a set of representatives for $\{f\restr{P}\mid P\in \partitn{\bar{y}}\}$
by \cref{FGconsequences}. Set $w=\gl_{g\in M} \omega H_g$. We claim that for all clopen set $U\ni \bar{y}$ we have
%Writing $(M_g)_{g\in M}$ \ynote*{unnecessary by \cref{cor:wedgeSets}}{for any injective enumeration of $\set{M_g}[g\in M]$}, \ynote*{for all clopen $U\ni \bar{y}$ in fact, do we use it?}{we claim that}
\[
f\leq w \glbin \fctwedge((M_g)_{ g\in M}\mid D) \leq f\corestr{U},
\]
which concludes the proof, once we prove the existence of the two reductions.

\smallskip

\textit{Right reduction}: $ w \glbin \fctwedge((M_g)_{ g\in M}\mid D) \leq f\corestr{U}$.
Let $W=\bigcup_{g\in M}W_g$ and note that by \cref{Intertwinereductionsforomegacentered} we have $w\leq f\corestr{W}$.
For each $g\in M$, we choose some $P_g\in\partit_M$ such that $f\restr{P_g}\equiv g$. We define $x_g$ to be any center for $f\restr{P_g}$.
Since $W$ is clopen and $\bar{y}\notin W$, we can pick a clopen $V\ni \bar{y}$ disjoint from $W$ and included in $U$.

We actually show that $\fctwedge((M_g)_{ g\in M}\mid D) \leq f\corestr{V}$ by verifying
the two conditions in the premice of \cref{DisjointificationLemma} for $\bar{y}$ and the points $x_g$.  

For the first condition, we need to show that for all $g\in M$, for every open $U'\ni x_g$ and all $h\in M_g$
there exists $(\sigma,\tau)$ reducing $h$ to $f$ with $\im(\sigma)\subseteq U'$ and $\overline{\im(f\sigma)}\notni \bar{y}$.
%\ynote*{it's not the first time we need to make this argument, maybe it could be stated in the section about centered functions}
{Since $x_g\in P_g\cap U'$ is a center for $f\restr{P_g}\equiv g$, there is a reduction $(\sigma,\tau)$ from $g$ to $f\restr{P_g\cap U'}$ by \cref{Centerinvariance}~\cref{Centerinvariance1}.
But the \cocenter{} of $f\restr{P_g\cap U'}$ is $\bar{y}$ so by \cref{Rigidityofthecocenter} we obtain that the desired reduction by restricting $(\sigma,\tau)$.  %$\overline{\im(f\sigma)}\notni \bar{y}$.
}
 
To verify the second condition, we need to show that
for every open $U'\ni \bar{y}$, there exists $(\sigma,\tau)$ reducing $\gl D$ to $f$ with $\im(f\sigma)\subseteq U'$ and $\overline{\im(f\sigma)}\notni \bar{y}$.
It is enough to find such a continuous reduction from $g$ to $f\corestr{U'}$ for every $g\in D$, so let $g\in D$.
Because $f$ is strongly solvable in $U$, there exists $P\in\partit_D$ such that $f\restr{P}\equiv g$ and $y_P\in U'$.
As $y_P\neq \bar{y}$, we can pick a clopen neighbourhood $V'\subseteq U'$ of $y_P$ that does not include $\bar{y}$.
Since $f\restr{P}$ is centered with \cocenter{} $y_P$ we have $g\leq f\corestr{V'}$, which gives the desired reduction. 

Altogether we obtain by \cref{Gluingasupperbound}:
\[
w \glbin \fctwedge((M_g)_{ g\in M}\mid D) \leq f\corestr{W}\glbin f\corestr{V} \leq f\corestr{U}.
\]

\smallskip
\textit{Left reduction:} $f\leq w \glbin \fctwedge((M_g)_{ g\in M}\mid D)$.
Let $A^0=\bigsqcup_{g\in M} A^0_g$, note that $A^0$ is clopen (since $M$ is finite) and $f\restr{A^0}\leq w$ by \cref{Gluingasupperbound}.
We show that for $A^1=A\setminus A^0$ we have $f\restr{A^1}\leq \fctwedge((M_g)_{ g\in M}\mid D)$ by verifying the hypotheses of \cref{Wedgeasupperbound}.
For all $g\in M$ we have $f\restr{A^1_g}\equiv g=\pgl M_g$ and the \cocenter{} of $f\restr{A^1_g}$ is $\bar{y}$.
By \cref{Rigidityofthecocenter} this implies that the sequence of rays $(\ray{(f\restr{A^1_g})}{\bar{y},j})_j$ is reducible by finite pieces to $\iw{M_g}$.

In case there is $P\in \partit_M$ with $\CB(f\restr{P})= \lambda+1$, let $A^1_0=\bigcup  (\partit_M\setminus \partit'_M)$.
Notice that $f\restr{A^1_0}$ is a simple function of rank $\lambda+1$ and distinguished point $\bar{y}$ by \cref{Simpleiffcoincidenceofcocenters}. So for all $j\in\N$ we have $\CB(\ray{(f\restr{A^1_0})}{\bar{y},j})\leq\lambda$ and by \cref{JSLgeneralstructure} we get $\ray{(f\restr{A^1_0})}{\bar{y},j}\leq \Maximalfct{\lambda}$.
Recall that any $g\in M$ has $\CB$-rank at least $\lambda+2$, so for some $h\in M_g$ we have $\CB(h)\geq\lambda+1$
and in turn $\Maximalfct{\lambda}\leq h$ by \cref{JSLgeneralstructure} again.
This means that we can safely replace $A^1_g$ by $A^1_g \sqcup A^1_0$ for some $g\in M$, since $(\ray{(f\restr{A^1_g})}{\bar{y},j})_j$ still reduces by pieces to $\iw{M_g}$.

If $\partitn{\bar{y}}$ is empty and so $D=\emptyset$, then $(A^1_g)_g\in M$ partitions $A^1$ and so $f\restr{A^1}\leq \fctwedge((M_g)_{ g\in M}\mid \emptyset)$ by \cref{Wedgeasupperbound}.
\smallskip

Otherwise, we have $A^1=(\bigsqcup_{g\in M} A^1_g)\sqcup A^D$ with $A^D=\bigcup \{P\mid P\in \partitn{\bar{y}}\}$, so we need to partition $A^D$ into an infinite sequence of clopen sets $(A^D_n)_{n\in\N}$
that is reducible by pieces to $\iw{D}$ and such that $f(A^D_n) \to \bar{y}$.

Fix any enumeration $(y_n)_{n\geq 1}$ of $Y'$ and set $A_n=\bigcup \{ P\in \partit \mid y_p=y_n\}$ for all $n\geq 1$.
For all $n\geq 1$ let $k_n\in\N$ be such that $y_n\in \ray{B}{\bar{y},k_n}$ and define $A^D_n=A_n\cap f^{-1}( \ray{B}{\bar{y},k_n})$ and $A^0_n=A_n\setminus A^D_n$. Finally let $A^D_0=\bigcup_{n\geq 0} A^0_n$. %Fix any enumeration $(y_n)_{n\geq 1}$ of $Y$ and set $A^D_n=A^D_{y_n}$ for all $n\geq 1$.
We have seen that strong solvability implies $y_n\to \bar{y}$, so $k_{n}\to\infty$, and in turn $f(A^D_n)\to \bar{y}$.

To conclude using \cref{Wedgeasupperbound}, it only remains to prove that $(A^D_n)_{n\in\N}$ is reducible by finite pieces to $\iw{D}$, or equivalently that $f\restr{A^D_n}\leq\FinGl{D}$ for all $n\in\N$.
By our definition, we need to distinguish two cases.

\textit{Case $n>0$.} We actually show the stronger relation $f\restr{A_n}\leq 2(\gl D)$.
We partition $A_n=\bigsqcup_{g\in D} A^g_n$ where $A^g_n=\bigcup\{P\in\partit \mid y_P=y_n \text{ and } f\restr{P}\equiv g\}$.
For all $g\in D$, $f\restr{A^g_n}$ is pseudo-centered, so by \cref{VerticalTheorem} we have $f\restr{A^g_n}\leq 2g$.
Therefore, by \cref{Gluingasupperbound},
\[ f\restr{A_n}\leq 2 (\gl D).\]


\textit{Case $n=0$.} We finally prove that $f\restr{A^D_0} \leq \FinGl{D}$. Note that $A^D_0=\bigsqcup_{P\in \partit_D}P\cap A^D_0$. Let $P\in \partitn{\bar{y}}$, $k$ be such that $y_P\in \ray{B}{\bar{y},k}$ and $g\in D$ with $f\restr{P}\equiv g$. Note that by definition $h:=f\restr{P\cap A^D_0}=f\restr{P}\corestr{B\setminus \ray{B}{\bar{y},k}}$.{If $g\equiv \Minimalfct{\lambda+1}$, then as $y_P$ is the distinguished point of $f\restr{P}$ we have $\CB(h)< \lambda$ }and so in particular $h\leq \Maximalfct{\lambda}$. Otherwise, $g\equiv \pgl D_g$ for some subset $D_g\subseteq  \centered{\alpha+1}\cup \omega \{\centered{\alpha+1}\}$ and so $h\leq \FinGl{D_g}$ by \cref{ResidualCorestrictionOfCentered}. Since $D$ is not empty and for any $g\in D$ we have $\Maximalfct{\lambda}\leq g$, we obtain:
\[
f\restr{A^D_0}\leq \gl_{P\in\partit_D} f\restr{P\cap A^D_y} \leq \gl_{g\in D}\omega D_g \leq  \gl_{g\in D}\pgl D_g \equiv \gl D 
\]
where we used that for each $g\in D$ we have $\omega D_g\leq \pgl D_g=g$ by \cref{GluinglowerthanPgluing}.
\end{proof}


\section{Solvable functions}
We finally introduce the weaker notion of \emph{solvable functions} by dropping the first requirement in the definition of strongly solvable functions.
We show that we can reduce the general case to the case of solvable functions and then we show how solvable functions can be expressed in terms of generators.  

\begin{definition}
We say that the function $f:A\to B$ is \emph{solvable} {with} a fine $c$-partition $\partit$ if for some $y\in Y_\partit$, for all $P\in\partitn{y}$ and all clopen $V\ni y$ there exists $Q\in \partitn{y}$ with $y_{Q}\in V$ and $f\restr{P}\leq f\restr{Q}$. \index[IdxNotion]{solvable function}\index[IdxNotion]{function!solvable}

\end{definition}
Note that if $f$ is solvable with $\partit$, then either $Y_\partit$ is a singleton or it is infinite.
If $\partit$ is a fine $c$-partition for a function $f:A\to B$ and $U$ is a clopen subset of $B$,
we define the corestriction $\partit\corestr{U}=\{P\in \partit \mid y_P\in U\}$ and we let $A^U_\partit=\bigcup \partit\corestr{U}$.\index[IdxSymb]{P@{$\partit\corestr{U}$}}\index[IdxSymb]{A@{$A^U_\partit$}}
Note that $\partit\corestr{U}$ is a fine $c$-partition of $f\restr{A^U_\partit}$, because any $\partit\corestr{U}$-lump is a $\partit$-lump.


\begin{theorem}\label{SolvableDecomposition}
For $\alpha<\omega_1$ assume $\FG(<\alpha+2)$ and let $\partit$ be a fine $c$-partition of $f:A \to B$ in $\sC_{\alpha+2}$.
Then there exists a countable family $\mathcal{U}$ of pairwise disjoint clopen subsets of $B$ such that:
\begin{enumerate}
\item $Y_\partit \subseteq \bigcup \mathcal{U}$, and 
\item for all $U\in \mathcal{U}$, the function $f\restr{A^U_\partit}:A^U_\partit\to B$ is solvable with $\partit\corestr{U}$.
\end{enumerate}
\end{theorem}
\begin{proof}
We start by observing that for every $y\in Y_\partit$ and every small enough neighborhood $U\ni y$ the function $f\restr{A^U_\partit}:A^U_\partit\to B$ is solvable in $\partit\corestr{U}$. 
To see this, fix a basis of clopen neighborhoods $(U_n)_{n\in \N}$ at $y$ and define:
\[
D_n=\set{g\in \centered{\alpha+2}}[\exists P\in\partitn{y} \text{ and } y_P\in U_n \text{ and } f\restr{P}\equiv g]
\]
Note that $m<n$ implies $D_m\supseteq D_n$, therefore as $\centered{\alpha+2}$ is finite this sequence must be eventually constant. So let $M\in\N$ be such that $D_M=D_n$ for all $n\geq M$ and let $U=U_M$. 
Clearly $\partit\corestr{U}$ is a fine $c$-partition of $f\restr{A^U_\partit}$. Suppose that $P\in \partit\corestr{U}$ with $y_P\neq y$ and let $V$ be a clopen neighborhood of $y$. As we assume $\FG(<\alpha+2)$, there exists $g\in \centered{\alpha+2}$ with $g\equiv f\restr{P}$ by \cref{FGconsequences}, so $g\in D_M$. Moreover there exists $n>M$ with $U_n\subseteq V$, hence as $D_n=D_M$, it follows that there exists $Q\in\partit\corestr{U}$ with $y_Q\in U_n\subseteq V$ and $f\restr{Q}\equiv f\restr{P}$.

Now fix an enumeration $(y_n)_n$ of $Y_\partit$. We define by induction a sequence $(U_k)_k$ such that the sets $U_k$ are pairwise disjoint clopen sets of $B$,
they cover $Y_\partit$, and for all $k$ the restriction $f\restr{A^{U_k}_\partit}:A^{U_k}_\partit\to B$ is solvable in $\partit\corestr{U_k}$.
Let $n_0=0$ and use the observation to choose a neighborhood $U_0$ of $y_0$ such that $f\restr{A^{U_0}_\partit}:A^{U_0}_\partit\to B$ is solvable in $\partit\corestr{U_0}$. Assume that $(U_l)_{l<k}$ is defined with $U_l$ clopen pairwise disjoint and $f\restr{A^{U_l}_\partit}:A^{U_l}_\partit\to B$ is solvable in $\partit\corestr{U_l}$ for all $l<k$. If $Y_P\subseteq \bigcup_{l<k}U_l$ we are done, otherwise let $n_k=\min \{n\mid y_n\notin  \bigcup_{l<k}U_l\}$
and using the observation again choose a small enough neighbourhood $U_k$ disjoint from $\bigcup_{l<k}U_l$ and such that $f\restr{A^{U_k}_\partit}:A^{U_k}_\partit\to B$ is solvable in $\partit\corestr{U_k}$. 
\end{proof}

\begin{remark}
If $f$ is strongly solvable at $y$ then we saw that $f\leq f\corestr{U}$ for every neighborhood $U$ of $y$.
In contrast, when $f$ is solvable with $\partit$ we will show that $f\leq f\corestr{U}$ whenever $U$ is a clopen set such that $Y_\partit\subseteq U$.
Note that it would be possible to strengthen the notion of solvable functions to a notion of \emph{solvable function at $y$} which further requires that
for any sequence $(P_n)_{n\in \N}$ in $\partitn{y}$ with pairwise distinct $y_{P_n}$ and every clopen set $V\ni y$
there exists a sequence $(Q_n)_{n\in \N}$ in $\partitn{y}$ with pairwise distinct $y_{Q_n}$ in $V$ such that $f\restr{P_n}\leq f\restr{Q_n}$.
This would ensure that $f\leq f\corestr{V}$ for any clopen neighborhood of $y$.
While this may seem desirable and a result analogous to \cref{SolvableDecomposition} can be established,
it seems to introduce a complexity which did not appear necessary for our proof.
\end{remark}

%At this point, our last task consists of showing
We now want to show the following statement for all $\alpha<\omega_1$:
\begin{descThm}
 \item[$\text{S}(\alpha)$] $\FG(\leq \alpha)$ \index[IdxSymb]{S@{$\text{S}(\alpha)$}}implies that for every $f\in \sC_{\alpha+1}$ if $f$ solvable with $\partit$, then there exists $g\in \generator{\alpha+1}$ with $f\leq g$ and $g\leq f\corestr{U}$ for every clopen $U\supseteq Y_\partit$.
  \end{descThm}
In \cref{FGatdoublesuccessors}, we prove that this is enough to obtain $\FG(\alpha+2)$ from $\FG(<\alpha+2)$, thus completing the inductive proof of \cref{PreciseStructureThm}.

This latter proof notably relies on a combined application of \cref{ExistenceFinePartitions,SolvableDecomposition}
to functions in $\sC_{\alpha+2}$ (for $\alpha=\lambda+n$ with $\lambda$ limit or null) to first get a fine $c$-partition, which is then partitioned to get solvable functions. %and then a solvable decomposition.
This process may yield some solvable functions in $\sC_{\lambda+1}$, and
while we have established $\FG(\lambda+1)$ this does not directly imply $\text{S}(\lambda)$\footnote{Nor does $\text{S}(\lambda)$ directly imply $\FG(\lambda+1)$,
which makes \cref{FGatsuccessoroflimit} still necessary.}.
So we now prove $\text{S}(\lambda)$ for $\lambda$ limit or null, before showing $\text{S}(\alpha)$ in the successor case.

%This proof notably relies on \cref{ExistenceFinePartitions}, which states that $\FG(<\alpha+2)$ ensures the existence of fine $c$-partition for functions $f\in \sC_{\alpha+2}$.
%However, when $\alpha=\lambda+n$ for $\lambda$ limit or null, some functions in $\sC_{\lambda+1}$ are actually solvable and they may arise from the application of \cref{SolvableDecomposition} to functions in $\sC_{\alpha+2}$. 
%While we have established that $\FG(\lambda+1)$ for $\lambda$ limit or null, note that this does not directly imply $\text{S}(\lambda)$. We therefore now prove $\text{S}(\lambda)$ for $\lambda$ limit or null before we show $\text{S}(\alpha)$ for successor ordinals.


\begin{proposition}\label{solvablelambda+1}
Let $\lambda<\omega_1$ be limit or null and assume $\FG(\leq \lambda)$. Suppose that $f:A\to B$ in $\sC_{\lambda+1}$ is solvable with $\partit$.

Then there exists a finite gluing $g$ of functions in $\generator{\lambda+1}$ such that $f\leq g$ and $g\leq f\corestr{U}$ for every clopen $U\supseteq Y_\partit$.
\end{proposition}

\begin{proof}
Let $\partit$ be a fine $c$-partition of $f$ and $y\in Y_\partit$ witnessing that $f$ is solvable and let $U\supseteq Y_\partit$ be clopen in $B$.
We first deal with the case $\lambda=0$.
As $f$ is in this case locally constant, we know that $f\equiv |\im f| \Minimalfct{1}$ from \cref{LocallyConstantFunctions}.
Since $f$ is solvable with $\partit$, by centeredness $f\restr{P}$ is constant for all $P\in \partit$, and so $\im f= Y_\partit\subseteq U$. So $f=f\corestr{U}$
and the result follows\footnote{Note that since $f$ is solvable, it actually follows that $\im f$ is either infinite or equal to $\set{y}$,
hence $f\equiv \omega \Minimalfct{1}$ or $f\equiv \Minimalfct{1}$.}.

Now assume that $\lambda$ is limit, recall that $\FG(\leq \lambda)$ ensures that $\FG(\lambda+1)$ holds by \cref{FGconsequences}, and using 
that $\partit$ is fine with \cref{cor:CenteredSucessor},
for every $P\in\partit$ we have $f\restr{P}\equiv \Minimalfct{\lambda+1}$ or $f\restr{P}\equiv  \pgl\Maximalfct{\lambda}$. 

\smallskip

First case:  $y_P=y$ for all $P\in \partit$, so $f$ is simple by \cref{Simpleiffcoincidenceofcocenters} and so $f\leq \pgl \Maximalfct{\lambda}$ by \cref{Maxfunctions}.
If $f\restr{P}\equiv  \pgl\Maximalfct{\lambda}$ for some $P\in \partit$, then $\pgl \Maximalfct{\lambda}\leq (f\restr{P})\corestr{V}\leq f\corestr{V}$ for every clopen neighbourhood $V$ of $y$ by centeredness of $f\restr{P}$. Otherwise $f\restr{P}\equiv\Minimalfct{\lambda+1}$ for all $P\in \partit$. If $\CB(\ray{f}{y,n})<\lambda$ for all $n\in \N$, then $f\leq \pgl_n \ray{f}{y,n}\leq \Minimalfct{\lambda+1}$ by \cref{Pgluingofraysasupperbound,ConsequencesGeneralStructureThm} and so we are done in this case too. Finally, assume that for some $n\in \N$ we have $\CB(\ray{f}{y,n})=\lambda$ and so $\Maximalfct{\lambda}\leq \ray{f}{y,n}$ by the General Structure \cref{JSLgeneralstructure}. As $\partit$ is a fine partition, $f_{y,\Minimalfct{\lambda+1}}=f$ is not a lump and so we can find a large enough $N\in \N$ such that $\Maximalfct{\lambda}\leq  \ray{f}{y,N}$ and $\ray{B}{y, N}\subseteq U$. Therefore for $V=U \setminus \ray{B}{y,N}$ we have  $\Minimalfct{\lambda+1}\glbin\Maximalfct{\lambda}\leq f\corestr{V} \glbin f\corestr{\ray{B}{y,N}}\leq f\corestr {U}$, as desired\footnote{Such a function $f\equiv \Minimalfct{\lambda+1}\glbin\Maximalfct{\lambda}$ can for example be constructed using the infinitary wedge (cf. \cref{rem:infinitewedge}) where each column consists of a sequence of functions $(f_{i,j})_{j\in\N}$ whose $\CB$-rank is cofinal in $\lambda$. While $\Minimalfct{\lambda+1}\glbin\Maximalfct{\lambda}$ does not admit a fine $c$-partition, the function $f$ clearly does.}.

\smallskip

Second case: there exists $P\in \partit$ with $y_P\neq y$. This means that $Y_\partit$ is actually infinite since $f$ is solvable.
We partition $\partit=\partit_0\sqcup\partit_1$ with $\partit_1=\set{P\in\partit}[f\restr{P}\equiv\pgl\Maximalfct{\lambda}]$ and set $Y_i=\set{y_P}[P\in\partit_i]$, $i=0,1$.
If $Y_1$ is empty then $Y=Y_0$ is infinite and $f\leq \omega\Minimalfct{\lambda+1}$ by \cref{Gluingasupperbound_cor}. Since $Y_0\subseteq U$, we get $\omega\Minimalfct{\lambda+1}\leq f\corestr{U}$ by \cref{Intertwinereductionsforomegacentered}, so we are done. If $Y_1$ is infinite, then $\Maximalfct{\lambda+1}=\omega \pgl \Maximalfct{\lambda} \leq f$ by \cref{Gluingasupperbound_cor} and $\Maximalfct{\lambda+1} \leq f\corestr{U}$ by \cref{Intertwinereductionsforomegacentered}, as desired. 

Suppose now that $Y_1$ is finite and nonempty. Then we claim that $Y_1=\set{y}$.  

Towards a contradiction, assume that $y_P\neq y$ for some $P\in \partit_1$. Then $f\restr{P}\equiv \pgl\Maximalfct{\lambda}$, but since $Y_1$ is finite we can find a small enough clopen neighbourhood $V$ of $y$ such that $V\cap Y_1=\emptyset$. But as $f$ is solvable, there exists $Q\in\partitn{y}$ with $y_Q\in V$ and $f\restr{P}\leq f\restr{Q}$. Since necessarily $f\restr{Q}\equiv \Minimalfct{\lambda+1}$, we get $\pgl\Maximalfct{\lambda}\leq \Minimalfct{\lambda+1}$, in contradiction with \cref{cor:CenteredSucessor}.
So we have $Y_1=\set{y}$ and $Y_0$ is infinite. We distinguish two subcases.
\begin{enumerate}
\item Suppose that for every clopen $V\ni y$ the set $Y_0\setminus V$ is finite. So $Y_1$ converges to $y$ and therefore \cref{Diagonalforlambda+1} allows to conclude.
\item Otherwise, $Y_0\cap W$ is infinite for some clopen $W\notni y$ and set $V=U\setminus W$. Let $f_i=\bigsqcup_{P\in \partit_i}f\restr{P}$ for $i=0,1$. By the first case, we have $f_1\leq \pgl \Maximalfct{\lambda}\leq f_1\corestr{V}$ and note that $f_1\corestr{V}\leq f\corestr{V}$. By \cref{Gluingasupperbound_cor,Intertwinereductionsforomegacentered}, we get $f_0\leq \omega\Minimalfct{\lambda+1}\leq f\corestr{W}$. Therefore
\[
f=f_1\sqcup f_0 \leq \pgl \Maximalfct{\lambda} \glbin \omega\Minimalfct{\lambda+1} \leq f\corestr{V}\glbin f\corestr{W}\leq f\corestr{U},
\]
as desired.\qedhere
\end{enumerate}
\end{proof}


\index[IdxNotion]{Precise Structure!solvable functions}
\index[IdxNotion]{finitely generated set of functions}
\index[IdxNotion]{generators}
\begin{theorem}\label{FiniteGenerationForSolvable}
Assume $\FG(\leq\alpha+1)$ for $\alpha<\omega_1$. Let $f:A \to B$ in $\sC_{\alpha+2}$ be solvable with $\partit$.
Then there exists $g\in \FinGl{\generator{\alpha+2}}$ such that $f\leq g$ and $g\leq f\corestr{U}$, so in particular $f\equiv g \equiv f\corestr{U}$, for every clopen $U\supseteq Y_\partit$.
\end{theorem}
\begin{proof}
Let $y\in Y_\partit$ witness that $f$ is solvable with $\partit$ and suppose that $U\supseteq Y_\partit$ is clopen in $B$.

By \cref{FGconsequences}, there exists a finite set $G\subseteq \centered{\alpha+2}$ of representatives for $\{f\restr{P}\mid P\in \partitn{y}\}$.
Let us say that $g\in G$ is \emph{reducible infinitely often off} $y$ if there exists a clopen neighborhood $V$ of $y$ in $B$
such that the set $V_g=\set{y_P}[P\in\partit \text{ and }y_P\notin V \text{ and }g\leq f\restr{P}]$ is infinite. We define
\[
H=\set{g\in G}[\text{$g$ is reducible infinitely often off $y$}]
\]
Note that $H$ is a $\leq$-downward closed subset of $G$ and set $D=G\setminus H$.
We further define $\partit_H=\{P\in\partitn{y}\mid \exists h\in H\,  f\restr{P}\equiv h\}$ and $\partit_M=\partit\setminus \partit_H$. 

We claim that the function $f_M=\bigsqcup_{P\in \partit_M}f\restr{P}$ together with $\partit_M$ is strongly solvable at $y$.
To this end, suppose that $W$ is a clopen neighborhood of $y$.
To check the first condition, suppose towards a contradiction that the set $\set{y_P}[P\in\partit_M \text{ and } y_P\notin W]$ is infinite.
Since $D$ is finite, the pigeonhole principle ensures that some $g\in D$ reduces infinitely often off $y$, a contradiction with the definition of $D$.
For the second condition, let $P\in\partit_M$ with $y_P\neq y$.
As $f$ is solvable with $\partit$ there exists $Q\in \partitn{y}$ such that $f\restr{P}\leq f\restr{Q}$.
Note that if $Q\in \partit_H$, then so is $P$ since $H$ is $\leq$-downward closed. Since $P\in \partit_M$, it follows that $Q\in \partit_M$ too, as desired.

Therefore $f_M$ is strongly solvable at $x$ and by \cref{DiagonalTheorem} there exists $g_M\in\FinGl{\generator{\alpha+2}}$ such that
\[
f_M \leq g_M \leq f_M\corestr{V}\leq f\corestr{V}.
\]
for every clopen neighborhood $V$ of $y$ in $B$. 

Finally, note that since $H$ is finite we can choose a single clopen neighborhood $V$ of $y$ in $B$ such that the sets $V_g$ are infinite for all $g\in H$.
It follows from \cref{Intertwinereductionsforomegacentered} that $\omega H\leq f\corestr{U\setminus V}$ for all clopen set $U\supseteq Y_\partit$.
Hence for $f_H=\bigsqcup_{P\in \partit_H}f\restr{P}$, the following holds by \cref{Gluingasupperbound}:
\[
f_H \leq \omega H \leq f\corestr{U\setminus V}.
\]

Altogether, for every clopen set $B\supseteq Y_\partit$ we have as desired: 
\[
f \leq f_H \glbin f_M \leq \omega H \glbin g_M \leq f\corestr{U\setminus V} \glbin f\corestr{V\cap U} \leq f\corestr{U}.\qedhere
\]
\end{proof}

\index[IdxNotion]{Precise Structure!double successors}
\index[IdxNotion]{finitely generated set of functions}
\index[IdxNotion]{generators}
\begin{theorem}\label{FGatdoublesuccessors}
For all $\alpha<\omega_1$, if $\FG(<\alpha+2)$ holds then so does $\FG(\leq\alpha+2)$.
\end{theorem}

\begin{proof}
Assume $\FG(<\alpha+2)$ and let $f:A\to B$ in $\sC_{\alpha+2}$, we prove that $f$ is equivalent to some function in $\FinGl{\generator{\alpha+2}}$.
By \cref{ExistenceFinePartitions} we can take a fine $c$-partition of $f$, and by \cref{SolvableDecomposition} there is a countable family $\mathcal{U}$
of pairwise disjoint clopen subsets of $B$ whose union covers $Y_\partit$ and such that $f_U=f\restr{A^U_\partit}$ is solvable for all $U\in \mathcal{U}$.
Now by \cref{solvablelambda+1,FiniteGenerationForSolvable}, for all $U\in \mathcal{U}$ there exists $g_U\in \FinGl{\generator{\alpha+2}}$
such that $f_U\leq g_U$ and $g_U\leq f_U\corestr{U}$, so in particular $f_U\equiv g_U \equiv f\corestr{U}$.
Putting everything together, by \cref{Gluingasupperbound_cor,Gluingaslowerbound2} we have
\[
f=\bigsqcup_{U\in\mathcal{U}}f_U\leq\gl_{U\in\mathcal{U}}f_U\leq\gl_{U\in\mathcal{U}}g_U\leq\gl_{U\in\mathcal{U}}{f_U}\corestr{U}\leq\gl_{U\in\mathcal{U}}f\corestr{U}\leq f.
\]
So $f\equiv\gl_{U\in\mathcal{U}}g_U$, which is as desired if $\mathcal{U}$ is finite. If $\mathcal{U}$ is infinite, we conclude using \cref{BasicsOnGenerators}~\cref{trickforfinalproof}.
\end{proof}

We can finally conclude the proofs of all our main results:

\index[IdxNotion]{Precise Structure!final proof}
\index[IdxNotion]{better-quasi-order (\bqo{})}
\begin{proof}[Proof of \cref{MainTheorem1,PreciseStructureThm,Introthm:LevelsAreFinitelyGenerated,Introthm:BQO,Introthm:BQOonScat,Introthm:BQO}]
As stated in \cref{FGconsequences}, we have already established $\FG(0)$, $\FG(\lambda)$ and $\FG(<\lambda+1)$ implies $\FG(\lambda+1)$, for every $\lambda$ limit.
\Cref{FGatdoublesuccessors} therefore completes the proof by induction of \cref{PreciseStructureThm}.
This in particular yields \cref{Introthm:LevelsAreFinitelyGenerated}, which in turn implies \cref{Introthm:BQO,Introthm:BQOonScat} by \cref{FGgivesBQO_2}.
This therefore concludes the proof of \cref{MainTheorem1}, since any \bqo{} is \wqo{}.
\end{proof}





